\scanpage{24}
dernière suit du cor. 1 du dit. On a
\[
0\unsure{-|\to}\unsure{\mathfrak N|N}\to A\to A'\to0
\]
d'où, comme $\mathfrak N$ est de dim $<n$ donc $H^n(\mathfrak N)=0$,
$H^n(A)\simeq H^n(A')$ donc $\Omega$ est aussi \unsure{fond|fond.} pour $A'$.

[N.B. $\mathfrak N$ est le plus grand idéal par lequel on
peut diviser \unsure{sans} que $\Omega$ cesse de \unsure{être|rester} \unsure{fond}]

\uline{Corollaire 1} L'image de $H^n(\Omega)\xrightarrow{\unsure{g|\eta|\mathfrak D}}\mathrm I$ est l'annulateur
de $\mathfrak N$, \inserted{\unsure{;|i.e.}\deleted{\unsure{et}} $\operatorname{Coker}g\simeq D(\mathfrak N)$.} \deleted{Donc $g$ est surjectif si et seulement
si \unsure{tous les} \unsure{les idéaux}} \inserted{Les conditions suivantes sont équiv.}

(i) $g$ est surjectif

(ii) $\mathfrak N=0$, i.e les idéaux premiers associés à $0$ de $A$
sont tous de \deleted{dim} corang $n$

(iii) $\Omega$ est un $A$-module fidèle.

En effet, la première assertion résulte du
fait que $\mathfrak N$ est le noyau de $A\to\operatorname{Hom}(\Omega,\Omega)$,
qui est transposé de $g$. D'où les équivalences.

\uline{Corollaire 2} $\operatorname{Ker}g\simeq D\bigl(\operatorname{Hom}(\Omega,\Omega)/A\,\operatorname{id}_{\Omega}\bigr)$.

Les conditions suivantes sont équivalentes

(i) $g$ est injectif \qquad (i bis) $H^n(\Omega)$ est comonogène

(ii) Tout end. de $\Omega$ est une homothétie

(ii bis) $\operatorname{Hom}(\Omega,\Omega)$ est monogène

La première formule est de dualité.
$(\mathrm i)\Rightarrow(\mathrm{i\,bis})$, \deleted{\unsure{$(\mathrm{ii})$}} $(\mathrm{i\,bis})\Longleftrightarrow(\unsure{\mathrm{ii}'|\mathrm{ii}}\,\mathrm{bis})$ et
% archive sheet number: 23

